\section{The underlying structures are a connected map}
A set tells us which objects are allowed. A relation specifies how pairs are connected or compared. An operation combines objects. A map translates one description into another. Additional structure determines which questions make sense: addition permits linear combinations, order permits extremal choices, distance permits approximation, and probability permits weighted averaging.

Ask what information a proof actually uses. Hall's theorem uses a finite adjacency relation and cardinalities, not coordinates of the vertices. The contraction theorem uses distances and completeness, not vector addition. Cauchy--Schwarz uses the algebra and positivity of a dot product. The sign-sequence argument translates pair products into row orthogonality and then into a counting restriction. Seeing what is \emph{not} used helps identify possible generalizations without guessing them indiscriminately.

A representation change is useful only when its correspondence is controlled. An inverse converts equal outputs to equal inputs. A quotient forgets distinctions shown to be irrelevant. A matrix represents a linear action in chosen coordinates. An indicator turns a counting question into a sum of zero-one values. These are different constructions, but each creates a language in which a desired step becomes easier to express and justify.
